\documentclass[11pt,a4paper]{article}

\usepackage[utf8]{inputenc}
\usepackage[T1]{fontenc}
\usepackage[spanish]{babel}
\usepackage{amsmath,amssymb,mathtools}
\usepackage{geometry}
\usepackage{enumitem}
\usepackage{array}
\usepackage{lmodern}
\usepackage{microtype}

\geometry{margin=2.4cm}
\setlength{\parindent}{0pt}
\setlength{\parskip}{0.5em}

\newcommand{\R}{\mathbb{R}}
\newcommand{\K}{\mathbb{K}}
\newcommand{\Ker}{\operatorname{Ker}}
\newcommand{\Ima}{\operatorname{Im}}
\newcommand{\rg}{\operatorname{rg}}
\newcommand{\Span}{\operatorname{L}}
\newcommand{\Mat}{\operatorname{M}}
\newcommand{\Id}{\operatorname{Id}}
\newcommand{\an}{\operatorname{an}}

\begin{document}

\begin{center}
	{\LARGE\bfseries Álgebra Lineal}\\[0.4em]
	{\Large Tema 2: Aplicaciones Lineales}
\end{center}

\tableofcontents
\newpage

\section{Definición y propiedades}

Sean $V$ y $W$ espacios vectoriales sobre el mismo cuerpo $\K$.

Una aplicación:

\[
f:V\longrightarrow W
\]

es lineal, o un homomorfismo, si:

\[
f(u+v)=f(u)+f(v)
\]

y:

\[
f(\alpha u)=\alpha f(u)
\]

para todo $u,v\in V$ y $\alpha\in\K$.

Equivalentemente:

\[
\boxed{
f(\alpha u+\beta v)
=
\alpha f(u)+\beta f(v)
}
\]

para todo $\alpha,\beta\in\K$ y $u,v\in V$.

\subsection{Propiedades}

Toda aplicación lineal verifica:

\[
f(0)=0
\]

\[
f(-u)=-f(u)
\]

Si $S\leq V$, entonces:

\[
f(S)\leq W
\]

Si:

\[
S=\Span\{u_1,\ldots,u_n\}
\]

entonces:

\[
f(S)
=
\Span\{f(u_1),\ldots,f(u_n)\}
\]

Si $T\leq W$, entonces:

\[
f^{-1}(T)\leq V
\]

Las aplicaciones lineales conservan la dependencia lineal:

\[
\{u_1,\ldots,u_n\}
\text{ l.d.}
\Longrightarrow
\{f(u_1),\ldots,f(u_n)\}
\text{ l.d.}
\]

En cambio, no conservan necesariamente la independencia lineal.

\section{Núcleo e imagen}

\subsection{Núcleo}

El núcleo de una aplicación lineal:

\[
f:V\longrightarrow W
\]

es:

\[
\boxed{
\Ker f
=
\{x\in V:f(x)=0_W\}
=
f^{-1}(0_W)
}
\]

El núcleo es un subespacio vectorial de $V$.

Además:

\[
\boxed{
f\text{ es inyectiva}
\iff
\Ker f=\{0\}
}
\]

\subsection{Imagen}

La imagen de $f$ es:

\[
\boxed{
\Ima f
=
\{y\in W:
\exists x\in V,\ f(x)=y
\}
}
\]

La imagen es un subespacio vectorial de $W$.

Si $B=\{v_1,\ldots,v_n\}$ es una base de $V$:

\[
\Ima f
=
\Span\{f(v_1),\ldots,f(v_n)\}
\]

\subsection{Rango}

El rango de una aplicación lineal es:

\[
\boxed{
\rg f=\dim(\Ima f)
}
\]

Si $V$ es de dimensión finita:

\[
\boxed{
\dim(\Ker f)+\dim(\Ima f)=\dim V
}
\]

Esta igualdad permite relacionar directamente núcleo, imagen e inyectividad.

\section{Clasificación de aplicaciones lineales}

Una aplicación lineal $f:V\to W$ puede clasificarse como:

\begin{itemize}
	\item \textbf{Endomorfismo}: si $V=W$.
	\item \textbf{Monomorfismo}: si es inyectiva.
	\item \textbf{Epimorfismo}: si es suprayectiva.
	\item \textbf{Isomorfismo}: si es biyectiva.
	\item \textbf{Automorfismo}: si es un endomorfismo biyectivo.
\end{itemize}

Se cumple:

\[
f\text{ inyectiva}
\iff
\Ker f=\{0\}
\]

\[
f\text{ suprayectiva}
\iff
\Ima f=W
\]

Si $V$ tiene dimensión finita:

\[
f\text{ inyectiva}
\iff
\dim V=\dim(\Ima f)
\]

Si $V$ y $W$ tienen la misma dimensión finita:

\[
f\text{ inyectiva}
\iff
f\text{ suprayectiva}
\iff
f\text{ biyectiva}
\]

\subsection{Isomorfismos}

Dos espacios vectoriales de dimensión finita son isomorfos si y sólo si tienen la misma dimensión.

Si:

\[
f:V\longrightarrow W
\]

es un isomorfismo, entonces:

\[
f^{-1}:W\longrightarrow V
\]

también es un isomorfismo.

\section{Espacio vectorial de aplicaciones lineales}

El conjunto de todas las aplicaciones lineales de $U$ en $V$ se denota:

\[
L(U,V)
\]

La suma se define mediante:

\[
(f+g)(u)=f(u)+g(u)
\]

y el producto por escalares:

\[
(\alpha f)(u)=\alpha f(u)
\]

Con estas operaciones:

\[
\boxed{
L(U,V)
\text{ es un espacio vectorial}
}
\]

Si:

\[
\dim U=n,
\qquad
\dim V=m
\]

entonces:

\[
\boxed{
\dim L(U,V)=mn
}
\]

\subsection{Composición}

Si:

\[
f:U\longrightarrow V,
\qquad
h:V\longrightarrow W
\]

la composición:

\[
h\circ f:U\longrightarrow W
\]

viene dada por:

\[
(h\circ f)(u)=h(f(u))
\]

y también es lineal.

\section{Aplicaciones lineales y matrices}

Sean:

\[
B=\{e_1,\ldots,e_n\}
\]

una base de $V$ y:

\[
B'=\{u_1,\ldots,u_m\}
\]

una base de $W$.

Si:

\[
f(e_j)
=
a_{1j}u_1+\cdots+a_{mj}u_m
\]

la matriz asociada a $f$ en las bases $B$ y $B'$ es:

\[
\boxed{
M_{B,B'}(f)
=
\begin{pmatrix}
|&&|\\
[f(e_1)]_{B'}&\cdots&[f(e_n)]_{B'}\\
|&&|
\end{pmatrix}
}
\]

Tiene dimensión:

\[
m\times n
\]

y transforma coordenadas según:

\[
\boxed{
[f(x)]_{B'}
=
M_{B,B'}(f)[x]_B
}
\]

La columna $j$ contiene las coordenadas de $f(e_j)$ en la base $B'$.

Además:

\[
\boxed{
\rg M_{B,B'}(f)=\dim(\Ima f)
}
\]

Por tanto:

\[
f\text{ inyectiva}
\iff
\rg M_{B,B'}(f)=n
\]

\[
f\text{ suprayectiva}
\iff
\rg M_{B,B'}(f)=m
\]

y:

\[
f\text{ isomorfismo}
\iff
M_{B,B'}(f)
\text{ es cuadrada y regular}
\]

\subsection{Operaciones con matrices asociadas}

\[
M_{B,B'}(f+g)
=
M_{B,B'}(f)+M_{B,B'}(g)
\]

\[
M_{B,B'}(\alpha f)
=
\alpha M_{B,B'}(f)
\]

Si:

\[
f:U\to V,
\qquad
h:V\to W
\]

entonces:

\[
\boxed{
M_{B,B''}(h\circ f)
=
M_{B',B''}(h)
M_{B,B'}(f)
}
\]

\section{Cambio de base}

Sean $A,A'$ bases de $U$ y $B,B'$ bases de $V$.

La matriz de una misma aplicación lineal en las nuevas bases viene dada por:

\[
\boxed{
M_{A',B'}(f)
=
M_{B,B'}
M_{A,B}(f)
M_{A',A}
}
\]

donde:

\[
M_{A',A}
\]

transforma coordenadas de $A'$ a $A$, y:

\[
M_{B,B'}
\]

transforma coordenadas de $B$ a $B'$.

\subsection{Endomorfismos}

Si $f:V\to V$ es un endomorfismo y $B,B'$ son dos bases de $V$:

\[
M_{B'}(f)
=
M_{B,B'}M_B(f)M_{B',B}
\]

Como:

\[
M_{B,B'}=M_{B',B}^{-1}
\]

si llamamos:

\[
P=M_{B',B}
\]

entonces:

\[
\boxed{
M_{B'}(f)
=
P^{-1}M_B(f)P
}
\]

Las matrices de un mismo endomorfismo en bases distintas son matrices semejantes.

\section{Matrices equivalentes, congruentes y semejantes}

\subsection{Matrices equivalentes}

Dos matrices $A,B\in M_{m\times n}(\K)$ son equivalentes si existen matrices regulares $P$ y $Q$ tales que:

\[
B=Q^{-1}AP
\]

Representan una misma aplicación lineal en bases diferentes y tienen el mismo rango:

\[
\boxed{
\rg A=\rg B
}
\]

\subsection{Matrices congruentes}

Dos matrices cuadradas $A$ y $B$ son congruentes si existe una matriz regular $P$ tal que:

\[
\boxed{
A=P^tBP
}
\]

Las matrices congruentes tienen el mismo rango.

\subsection{Matrices semejantes}

Dos matrices cuadradas $A$ y $B$ son semejantes si existe una matriz regular $P$ tal que:

\[
\boxed{
B=P^{-1}AP
}
\]

Representan un mismo endomorfismo en bases distintas.

\section{Descomposición canónica}

Sea:

\[
f:U\longrightarrow V
\]

una aplicación lineal.

El primer teorema de isomorfía establece:

\[
\boxed{
U/\Ker f
\cong
\Ima f
}
\]

mediante:

\[
\widetilde f:
U/\Ker f
\longrightarrow
\Ima f
\]

\[
\widetilde f(u+\Ker f)=f(u)
\]

La aplicación $f$ puede descomponerse como:

\[
U
\xrightarrow{\pi}
U/\Ker f
\xrightarrow{\widetilde f}
\Ima f
\xrightarrow{i}
V
\]

donde:

\[
\pi(u)=u+\Ker f
\]

es la proyección canónica,

\[
\widetilde f
\]

es un isomorfismo, e:

\[
i:\Ima f\hookrightarrow V
\]

es la inclusión.

Por tanto:

\[
\boxed{
f=i\circ\widetilde f\circ\pi
}
\]

Toda aplicación lineal puede expresarse como composición de un epimorfismo, un isomorfismo y un monomorfismo.

\section{Espacio dual y formas lineales}

Una forma lineal sobre $V$ es una aplicación lineal:

\[
f:V\longrightarrow\K
\]

El espacio dual de $V$ es:

\[
\boxed{
V^*=L(V,\K)
}
\]

Si:

\[
B=\{v_1,\ldots,v_n\}
\]

es una base de $V$, la matriz de una forma lineal $f$ es una matriz fila:

\[
\boxed{
M_B(f)
=
\begin{pmatrix}
f(v_1)&\cdots&f(v_n)
\end{pmatrix}
}
\]

Si $\dim V=n$:

\[
\boxed{
\dim V^*=n
}
\]

\section{Base dual}

Sea:

\[
B=\{v_1,\ldots,v_n\}
\]

una base de $V$.

Su base dual:

\[
B^*=
\{\varphi_1,\ldots,\varphi_n\}
\]

se caracteriza por:

\[
\boxed{
\varphi_i(v_j)=\delta_{ij}
}
\]

donde:

\[
\delta_{ij}
=
\begin{cases}
1, & i=j\\
0, & i\neq j
\end{cases}
\]

Además, si:

\[
[x]_B=(x_1,\ldots,x_n)
\]

entonces:

\[
\boxed{
x_i=\varphi_i(x)
}
\]

Si los vectores de $B$ escritos por columnas forman la matriz $A$, las formas lineales de la base dual vienen dadas por las filas de:

\[
\boxed{
A^{-1}
}
\]

\section{Anulador de un subespacio}

Sea $S\leq V$.

El anulador de $S$ es:

\[
\boxed{
\an(S)
=
\{
f\in V^*:
f(v)=0
\quad
\forall v\in S
\}
}
\]

También se denomina subespacio ortogonal a $S$.

Se cumple:

\[
\an(S)\leq V^*
\]

y:

\[
\an(S)=\an(\Span(S))
\]

por lo que basta imponer que una forma lineal anule a un sistema generador de $S$.

\section{Aplicación lineal traspuesta}

Sea:

\[
f:V\longrightarrow V'
\]

una aplicación lineal.

Se define su traspuesta:

\[
f^t:(V')^*\longrightarrow V^*
\]

mediante:

\[
\boxed{
f^t(\varphi)
=
\varphi\circ f
}
\]

Si la matriz de $f$ en unas bases es:

\[
A=M_{B,B'}(f)
\]

entonces la matriz de la aplicación traspuesta en las bases duales correspondientes es:

\[
\boxed{
M_{(B')^*,B^*}(f^t)
=
A^t
}
\]

\end{document}
